% Glossary entries (must be defined in the preamble for no-index mode).
% First body mention of each term uses \gls{<label>}, which renders the term
% and hyperlinks back to the Glossary.
\newglossaryentry{dtu}{
  name={Differential Transcript Usage (DTU)},
  first={Differential Transcript Usage (DTU)},
  text={DTU},
  description={Differential transcript usage occurs when the relative expression of multiple transcripts arising from the same gene changes between different conditions \cite{young2024dtu}. It is independent of changes in the gene's total expression, so it is not captured by gene-level differential expression analysis and requires transcript-level resolution to detect.}
}

\newglossaryentry{deg}{
  name={differential gene expression (DGE)},
  first={differential gene expression (DGE)},
  text={DGE},
  description={The analysis of differences in a gene's total expression---the summed expression of all its isoforms---between groups of samples. It measures overall gene output and is blind to changes in the relative usage of individual isoforms, which are instead captured by differential transcript usage.}
}

\newglossaryentry{threePrimeEnd}{
  name={3' end},
  description={The terminal end of an RNA molecule at which transcription finishes; in mRNA, it is where the poly-A tail is added. In 3'-biased protocols, the capture chemistry is anchored at this end, so reads cover only the terminal fragment of each molecule rather than the whole transcript.}
}

\newglossaryentry{threePrime}{
  name={3'-biased sequencing},
  first={3'-biased sequencing},
  text={3'-biased},
  plural={3'-biased},
  description={A property of some scRNA-seq protocols (such as 10x Genomics Chromium) in which sequencing reads are derived only from the 3' terminus of each transcript. Because reading quality decreases by distance and because capture is anchored at the transcript's 3' end, coverage does not extend into the transcript body, so internal splice junctions are never observed \cite{regan2022practical}.}
}

\newglossaryentry{isoformSwitch}{
  name={isoform switching},
  description={A change between conditions in the relative usage of a gene's isoforms, such that the fraction of the gene's expression contributed by one isoform increases while that of another decreases. It reflects altered alternative splicing or transcription and is quantified by the change in each isoform's fraction of total gene expression \cite{vitting2017isoformswitch}.}
}

\newglossaryentry{sparsity}{
  name={sparsity},
  description={The very high proportion of zero counts in single-cell RNA-seq data.}
}

\newglossaryentry{dropout}{
  name={technical dropout},
  plural={technical dropouts},
  description={A zero count produced by a technical limitation of the assay---failed capture, reverse transcription, or sequencing---rather than by the true absence of the transcript. It is also described as a non-biological zero.}
}

% No extra period between a description and its page list.
\renewcommand*{\glspostdescription}{}

\newglossaryentry{countMatrix}{
  name={count matrix},
  plural={count matrices},
  description={A table whose rows index features (genes or transcripts) and whose columns index cells, with each entry recording how many molecules of that feature were detected in that cell.}
}

\newglossaryentry{isoform}{
  name={isoform},
  plural={isoforms},
  description={A distinct mature transcript variant produced from the same gene by alternative splicing.}
}

\newglossaryentry{scRnaSeq}{
  name={single-cell RNA sequencing (scRNA-seq)},
  first={Single-cell RNA sequencing (scRNA-seq)},
  text={scRNA-seq},
  description={RNA sequencing performed on thousands of individual cells separately, exposing per-cell expression.}
}

\newglossaryentry{bulkRnaSeq}{
  name={bulk RNA-seq},
  description={Sequencing of RNA pooled from many cells at once, which averages away cell-to-cell differences.}
}

\newglossaryentry{fullLength}{
  name={full-length sequencing},
  description={A sequencing approach that reads an RNA molecule from one end to the other, rather than only the fragment nearest one of its ends.}
}

\newglossaryentry{read}{
  name={read},
  plural={reads},
  description={A short nucleotide sequence produced by the sequencing machine, representing the raw data of RNA-seq.}
}

\newglossaryentry{dropletProtocols}{
  name={droplet-based protocols (droplet data)},
  first={droplet protocols},
  text={droplet protocols},
  description={High-throughput single-cell protocols that encapsulate each cell together with a barcoded bead inside a picolitre droplet, so that tens of thousands of cells are processed in one run. They are the dominant scRNA-seq format, and their capture chemistry reads only the tail of each transcript.}
}

\newglossaryentry{tagBias}{
  name={tag bias},
  first={tag bias},
  description={A bias in which each molecule is read only from one of its ends, so that coverage concentrates near that end while the interior of the transcript is never observed. It follows from a capture chemistry anchored at a single end, as in the tail-primed droplet protocols.}
}

\newglossaryentry{cellBarcode}{
  name={cell barcode},
  plural={cell barcodes},
  description={A short DNA tag carried by every read from a given cell, identifying its cell of origin. It allows the reads of one pooled sequencing run to be regrouped into per-cell expression profiles.}
}

\newglossaryentry{umi}{
  name={Unique Molecular Identifier (UMI)},
  first={Unique Molecular Identifier (UMI)},
  text={UMI},
  plural={UMIs},
  description={A short random sequence attached to each captured RNA molecule before PCR amplification, so that amplified copies of one molecule are recognised and counted only once.}
}

\newglossaryentry{pcr}{
  name={PCR},
  first={PCR (polymerase chain reaction)},
  description={Polymerase chain reaction: the amplification step that copies each captured cDNA fragment many times over, so that the few molecules recovered from a cell are present in enough quantity for the sequencer to read \cite{saiki1985enzymatic}. Because copies outnumber original molecules, the UMI tag is what allows them to be collapsed back to a single count.}
}

\newglossaryentry{fastq}{
  name={FASTQ},
  description={The plain-text de facto standard output of Illumina sequencers, named by analogy to FASTA with a Q for quality. Files are almost always gzip-compressed, taking the \texttt{.fastq.gz} extension.}
}

\newglossaryentry{rna}{
  name={RNA (ribonucleic acid)},
  first={RNA (ribonucleic acid)},
  text={RNA},
  description={The broad class of nucleic acid molecules transcribed from DNA. It comprises messenger RNA (mRNA), which carries protein-coding instructions, and non-coding RNAs that perform structural or regulatory roles.}
}

\newglossaryentry{mrna}{
  name={messenger RNA (mRNA)},
  first={messenger RNA (mRNA)},
  text={mRNA},
  description={The protein-coding subtype of RNA that carries a gene's sequence to the ribosome.}
}

\newglossaryentry{rnaSeq}{
  name={RNA sequencing (RNA-seq)},
  first={RNA sequencing (RNA-seq)},
  text={RNA-seq},
  description={A technique that determines the sequence and abundance of RNA molecules in a sample; it is the basis of the single-cell protocols discussed in this thesis.}
}

\newglossaryentry{splicing}{
  name={alternative splicing},
  first={alternative splicing},
  text={alternative splicing},
  description={The process by which a gene's long primary RNA copy is shortened and rearranged during maturation. The copy contains coding blocks (exons) and non-coding blocks (introns); splicing removes the introns and joins the flanking exons at the resulting splice junctions. When exon selection varies between transcripts of the same gene, the result is multiple distinct isoforms.}
}

\newglossaryentry{mutualExclusivity}{
  name={mutual exclusivity},
  description={Two features that rarely appear together in the same cell (a negative association), as opposed to co-expression.}
}

\newglossaryentry{gdi}{
  name={Global Differentiation Index (GDI)},
  first={Global Differentiation Index (GDI)},
  text={GDI},
  description={A cluster-specificity measure used to identify marker genes and validate clustering quality.}
}

\newglossaryentry{coexpression}{
  name={co-expression},
  description={The tendency of two features (genes or transcripts) to be detected together across cells.}
}

\newglossaryentry{contingencyTable}{
  name={contingency table},
  description={A $2 \times 2$ table counting, across all cells, the four presence/absence combinations of two features; it tests whether they co-occur more or less than chance.}
}

\newglossaryentry{pseudoBulk}{
  name={pseudo-bulk aggregation},
  description={The practice of summing read counts across groups of single cells to simulate bulk RNA-seq depth, which recovers statistical power at the expense of single-cell heterogeneity.}
}

\newglossaryentry{orthogonalValidation}{
  name={orthogonal validation},
  description={Independent confirmation of a computational result by a different experimental technique.}
}

\newglossaryentry{transcript}{
  name={transcript},
  plural={transcripts},
  description={An RNA molecule copied from a gene.}
}

\newglossaryentry{unsplicedTranscript}{
  name={unspliced/spliced transcript},
  first={unspliced transcript},
  text={unspliced transcript},
  plural={unspliced transcripts},
  description={The first RNA copy of a gene, still containing its non-coding blocks (introns); it is also called the primary transcript or pre-mRNA. Its counterpart is the spliced transcript, the same molecule once those blocks have been removed, which is the mature form translated into protein.}
}

\newglossaryentry{spliceJunction}{
  name={splice junction},
  plural={splice junctions},
  description={The boundary created where two coding blocks (exons) are joined after the intervening non-coding block (intron) has been removed. Reads spanning a splice junction identify which blocks were joined together.}
}

\newglossaryentry{gene}{
  name={gene},
  plural={genes},
  description={A functional unit of the genome that specifies one or more RNA molecules.}
}

\newglossaryentry{protein}{
  name={protein},
  plural={proteins},
  description={A molecule built from a chain of amino acids that performs the work of the cell, including catalysing reactions, transmitting signals, and forming structural components.}
}

\newglossaryentry{genome}{
  name={genome},
  description={The complete set of genetic material of an organism, written in DNA and inherited from one generation to the next. It is organised into genes, the functional units whose expression is measured.}
}
% Suppress the package-generated heading; the chapter provides its own.
\renewcommand{\glossarysection}[2][]{}
% Print cue: mark glossary terms at first use so paper readers spot them.
\renewcommand{\glstextformat}[1]{\textit{#1}}
